%%%%%%%%%%%%%%%%%%%%%%%%%%%%%%%%%%%%%%%%%%%%%%%%%%%%%%%%%%%%
% 4.3 DIFFERENTIAL TOPOLOGY
% The Composition of Advanced Mathematics
%%%%%%%%%%%%%%%%%%%%%%%%%%%%%%%%%%%%%%%%%%%%%%%%%%%%%%%%%%%%

\section{Differential Topology}
\label{sec:differential-topology}

\prerequisite{
Point-Set Topology, Algebraic Topology, Multivariable Calculus,
Linear Algebra, and introductory Differential Geometry.
}

\learningobjectives{
By the end of this section, the reader should be able to:
\begin{itemize}
    \item explain the distinction between differential topology and differential geometry;
    \item define smooth manifolds and smooth maps;
    \item define immersions, submersions, and embeddings;
    \item identify regular and critical points;
    \item define regular and critical values;
    \item apply the Regular Value Theorem;
    \item state the Constant Rank Theorem;
    \item explain tangent maps and rank;
    \item describe submanifolds;
    \item explain transversality;
    \item state the Transversality Theorem conceptually;
    \item describe orientations of manifolds;
    \item define manifolds with boundary;
    \item explain degree and mod-\(2\) degree;
    \item state Sard's Theorem;
    \item describe vector fields and their zeros;
    \item state the Hairy Ball Theorem;
    \item explain the Poincar\'e--Hopf Theorem;
    \item describe isotopies;
    \item distinguish isotopy from homotopy;
    \item explain Morse functions and critical points;
    \item state the Morse Lemma;
    \item interpret Morse index;
    \item explain how topology changes at critical values;
    \item describe handle attachments;
    \item explain the role of cobordism;
    \item connect smooth structure with topological invariants.
\end{itemize}
}

%%%%%%%%%%%%%%%%%%%%%%%%%%%%%%%%%%%%%%%%%%%%%%%%%%%%%%%%%%%%
\subsection{What Is Differential Topology?}
\label{subsec:what-is-differential-topology}
%%%%%%%%%%%%%%%%%%%%%%%%%%%%%%%%%%%%%%%%%%%%%%%%%%%%%%%%%%%%

Differential topology studies topological spaces equipped with enough
smooth structure to use calculus.

Its central objects are smooth manifolds and smooth maps.

The subject lies between topology and differential geometry:

\[
\boxed{
\text{Topology}
\longleftrightarrow
\text{Differential Topology}
\longleftrightarrow
\text{Differential Geometry}.
}
\]

Topology asks which properties survive continuous deformation.

Differential topology asks which properties survive \emph{smooth}
deformation.

Differential geometry additionally studies geometric quantities such as

\[
\boxed{
\text{distance},
\quad
\text{angle},
\quad
\text{curvature},
\quad
\text{geodesics}.
}
\]

\begin{conceptbox}
Differential topology studies the global consequences of local smooth
structure without requiring a metric or curvature.
\end{conceptbox}

%%%%%%%%%%%%%%%%%%%%%%%%%%%%%%%%%%%%%%%%%%%%%%%%%%%%%%%%%%%%
\subsection{Smooth Manifolds}
\label{subsec:differential-topology-smooth-manifolds}
%%%%%%%%%%%%%%%%%%%%%%%%%%%%%%%%%%%%%%%%%%%%%%%%%%%%%%%%%%%%

A topological \(n\)-manifold locally resembles

\[
\R^n.
\]

To perform calculus, the coordinate charts must be compatible
differentiably.

\begin{definitionbox}{Smooth Manifold}
A \emph{smooth manifold} is a topological manifold equipped with an atlas
whose transition maps are infinitely differentiable.
\end{definitionbox}

If

\[
(U,\varphi)
\qquad\text{and}\qquad
(V,\psi)
\]

are overlapping charts, then the transition map

\[
\boxed{
\psi\circ\varphi^{-1}
}
\]

must be smooth wherever defined.

%%%%%%%%%%%%%%%%%%%%%%%%%%%%%%%%%%%%%%%%%%%%%%%%%%%%%%%%%%%%
\subsection{Smooth Maps}
\label{subsec:smooth-maps-differential-topology}
%%%%%%%%%%%%%%%%%%%%%%%%%%%%%%%%%%%%%%%%%%%%%%%%%%%%%%%%%%%%

Let \(M\) and \(N\) be smooth manifolds.

A map

\[
f:M\to N
\]

is smooth if its coordinate representations are smooth maps between
Euclidean spaces.

Locally,

\[
\boxed{
\psi\circ f\circ\varphi^{-1}
}
\]

must be infinitely differentiable.

%%%%%%%%%%%%%%%%%%%%%%%%%%%%%%%%%%%%%%%%%%%%%%%%%%%%%%%%%%%%
\subsection{Diffeomorphisms}
\label{subsec:diffeomorphisms-differential-topology}
%%%%%%%%%%%%%%%%%%%%%%%%%%%%%%%%%%%%%%%%%%%%%%%%%%%%%%%%%%%%

\begin{definitionbox}{Diffeomorphism}
A map

\[
f:M\to N
\]

is a \emph{diffeomorphism} if:

\begin{enumerate}
    \item \(f\) is bijective;
    \item \(f\) is smooth;
    \item \(f^{-1}\) is smooth.
\end{enumerate}
\end{definitionbox}

If such a map exists, we write

\[
\boxed{
M\cong_{\mathrm{diff}}N.
}
\]

Diffeomorphic manifolds are considered equivalent from the viewpoint of
smooth structure.

%%%%%%%%%%%%%%%%%%%%%%%%%%%%%%%%%%%%%%%%%%%%%%%%%%%%%%%%%%%%
\subsection{Homeomorphism Versus Diffeomorphism}
\label{subsec:homeomorphism-vs-diffeomorphism}
%%%%%%%%%%%%%%%%%%%%%%%%%%%%%%%%%%%%%%%%%%%%%%%%%%%%%%%%%%%%

Every diffeomorphism is a homeomorphism.

However, a homeomorphism need not be a diffeomorphism.

Therefore,

\[
\boxed{
\text{diffeomorphic}
\Longrightarrow
\text{homeomorphic},
}
\]

but the converse need not hold.

In sufficiently high dimensions, a single topological manifold may admit
distinct smooth structures.

This is one of the major discoveries of modern differential topology.

%%%%%%%%%%%%%%%%%%%%%%%%%%%%%%%%%%%%%%%%%%%%%%%%%%%%%%%%%%%%
\subsection{Tangent Spaces}
\label{subsec:differential-topology-tangent-spaces}
%%%%%%%%%%%%%%%%%%%%%%%%%%%%%%%%%%%%%%%%%%%%%%%%%%%%%%%%%%%%

At each point

\[
p\in M
\]

of an \(n\)-dimensional smooth manifold, there is an \(n\)-dimensional
vector space

\[
\boxed{
T_pM
}
\]

called the tangent space.

It represents the possible infinitesimal directions of motion through
\(p\).

%%%%%%%%%%%%%%%%%%%%%%%%%%%%%%%%%%%%%%%%%%%%%%%%%%%%%%%%%%%%
\subsection{Differential of a Smooth Map}
\label{subsec:differential-smooth-map}
%%%%%%%%%%%%%%%%%%%%%%%%%%%%%%%%%%%%%%%%%%%%%%%%%%%%%%%%%%%%

A smooth map

\[
f:M\to N
\]

induces a linear map

\[
\boxed{
df_p:T_pM\to T_{f(p)}N.
}
\]

This map is called the \emph{differential}, \emph{derivative}, or
\emph{pushforward} of \(f\) at \(p\).

In Euclidean coordinates, \(df_p\) is represented by the Jacobian matrix.

%%%%%%%%%%%%%%%%%%%%%%%%%%%%%%%%%%%%%%%%%%%%%%%%%%%%%%%%%%%%
\subsection{Rank}
\label{subsec:differential-topology-rank}
%%%%%%%%%%%%%%%%%%%%%%%%%%%%%%%%%%%%%%%%%%%%%%%%%%%%%%%%%%%%

The rank of \(f\) at \(p\) is

\[
\boxed{
\operatorname{rank}_p(f)
=
\operatorname{rank}(df_p).
}
\]

If

\[
f:\R^m\to\R^n,
\]

then this is the rank of the Jacobian matrix

\[
Df(p).
\]

Rank measures how many independent output directions are produced locally
by the map.

%%%%%%%%%%%%%%%%%%%%%%%%%%%%%%%%%%%%%%%%%%%%%%%%%%%%%%%%%%%%
\subsection{Immersions}
\label{subsec:immersions}
%%%%%%%%%%%%%%%%%%%%%%%%%%%%%%%%%%%%%%%%%%%%%%%%%%%%%%%%%%%%

\begin{definitionbox}{Immersion}
A smooth map

\[
f:M\to N
\]

is an \emph{immersion} if

\[
\boxed{
df_p:T_pM\to T_{f(p)}N
}
\]

is injective for every \(p\in M\).
\end{definitionbox}

If

\[
\dim M=m,
\]

then an immersion has

\[
\boxed{
\operatorname{rank}(df_p)=m.
}
\]

Thus an immersion locally places \(M\) smoothly inside \(N\), although
global self-intersections may occur.

%%%%%%%%%%%%%%%%%%%%%%%%%%%%%%%%%%%%%%%%%%%%%%%%%%%%%%%%%%%%
\subsection{Example of an Immersion with Self-Intersection}
\label{subsec:immersion-self-intersection}
%%%%%%%%%%%%%%%%%%%%%%%%%%%%%%%%%%%%%%%%%%%%%%%%%%%%%%%%%%%%

A smooth figure-eight curve in the plane may be parametrized so that its
velocity never vanishes.

Then

\[
f:S^1\to\R^2
\]

may be an immersion even though two distinct points of \(S^1\) have the same
image.

Therefore

\[
\boxed{
\text{immersion does not imply embedding}.
}
\]

%%%%%%%%%%%%%%%%%%%%%%%%%%%%%%%%%%%%%%%%%%%%%%%%%%%%%%%%%%%%
\subsection{Submersions}
\label{subsec:submersions}
%%%%%%%%%%%%%%%%%%%%%%%%%%%%%%%%%%%%%%%%%%%%%%%%%%%%%%%%%%%%

\begin{definitionbox}{Submersion}
A smooth map

\[
f:M\to N
\]

is a \emph{submersion} if

\[
\boxed{
df_p:T_pM\to T_{f(p)}N
}
\]

is surjective for every \(p\in M\).
\end{definitionbox}

If

\[
\dim N=n,
\]

then

\[
\boxed{
\operatorname{rank}(df_p)=n.
}
\]

%%%%%%%%%%%%%%%%%%%%%%%%%%%%%%%%%%%%%%%%%%%%%%%%%%%%%%%%%%%%
\subsection{Embeddings}
\label{subsec:differential-topology-embeddings}
%%%%%%%%%%%%%%%%%%%%%%%%%%%%%%%%%%%%%%%%%%%%%%%%%%%%%%%%%%%%

\begin{definitionbox}{Smooth Embedding}
A smooth map

\[
f:M\to N
\]

is a \emph{smooth embedding} if it is an immersion and a homeomorphism from
\(M\) onto its image

\[
f(M)
\]

with the subspace topology.
\end{definitionbox}

An embedding realizes \(M\) as a genuine smooth submanifold of \(N\).

%%%%%%%%%%%%%%%%%%%%%%%%%%%%%%%%%%%%%%%%%%%%%%%%%%%%%%%%%%%%
\subsection{Immersion Versus Embedding}
\label{subsec:immersion-vs-embedding}
%%%%%%%%%%%%%%%%%%%%%%%%%%%%%%%%%%%%%%%%%%%%%%%%%%%%%%%%%%%%

The distinction is global.

An immersion must behave correctly infinitesimally.

An embedding must additionally behave correctly globally.

Thus

\[
\boxed{
\text{embedding}
\Longrightarrow
\text{immersion},
}
\]

but not conversely.

%%%%%%%%%%%%%%%%%%%%%%%%%%%%%%%%%%%%%%%%%%%%%%%%%%%%%%%%%%%%
\subsection{Whitney Embedding Theorem}
\label{subsec:whitney-embedding}
%%%%%%%%%%%%%%%%%%%%%%%%%%%%%%%%%%%%%%%%%%%%%%%%%%%%%%%%%%%%

A foundational theorem states that abstract smooth manifolds can be realized
inside Euclidean spaces.

\begin{theorem}[Whitney Embedding Theorem]
Every smooth \(n\)-dimensional manifold satisfying the standard
second-countability assumptions admits a smooth embedding into

\[
\boxed{
\R^{2n}.
}
\]
\end{theorem}

Stronger refinements are possible in many cases.

The theorem shows that abstract manifolds can be studied concretely as
submanifolds of Euclidean space.

%%%%%%%%%%%%%%%%%%%%%%%%%%%%%%%%%%%%%%%%%%%%%%%%%%%%%%%%%%%%
\subsection{Whitney Immersion Theorem}
\label{subsec:whitney-immersion}
%%%%%%%%%%%%%%%%%%%%%%%%%%%%%%%%%%%%%%%%%%%%%%%%%%%%%%%%%%%%

A related theorem states that every smooth \(n\)-manifold admits an
immersion into

\[
\boxed{
\R^{2n-1}
}
\]

for \(n>1\), under the usual manifold hypotheses.

The embedding theorem is stronger because embeddings forbid global
self-intersections.

%%%%%%%%%%%%%%%%%%%%%%%%%%%%%%%%%%%%%%%%%%%%%%%%%%%%%%%%%%%%
\subsection{Critical Points}
\label{subsec:critical-points-differential-topology}
%%%%%%%%%%%%%%%%%%%%%%%%%%%%%%%%%%%%%%%%%%%%%%%%%%%%%%%%%%%%

Consider a smooth map

\[
f:M\to N.
\]

\begin{definitionbox}{Critical Point}
A point \(p\in M\) is a \emph{critical point} of \(f\) if

\[
\boxed{
df_p
}
\]

fails to have maximal possible rank.
\end{definitionbox}

For a real-valued function

\[
f:M\to\R,
\]

this means

\[
\boxed{
df_p=0.
}
\]

%%%%%%%%%%%%%%%%%%%%%%%%%%%%%%%%%%%%%%%%%%%%%%%%%%%%%%%%%%%%
\subsection{Regular Points}
\label{subsec:regular-points}
%%%%%%%%%%%%%%%%%%%%%%%%%%%%%%%%%%%%%%%%%%%%%%%%%%%%%%%%%%%%

A point that is not critical is called a \emph{regular point}.

For

\[
f:M\to\R,
\]

a point \(p\) is regular precisely when

\[
\boxed{
df_p\neq0.
}
\]

In Euclidean coordinates this corresponds to

\[
\boxed{
\nabla f(p)\neq0.
}
\]

%%%%%%%%%%%%%%%%%%%%%%%%%%%%%%%%%%%%%%%%%%%%%%%%%%%%%%%%%%%%
\subsection{Critical and Regular Values}
\label{subsec:critical-regular-values}
%%%%%%%%%%%%%%%%%%%%%%%%%%%%%%%%%%%%%%%%%%%%%%%%%%%%%%%%%%%%

\begin{definitionbox}{Critical Value}
A point \(y\in N\) is a \emph{critical value} of \(f\) if

\[
y=f(p)
\]

for some critical point \(p\).
\end{definitionbox}

\begin{definitionbox}{Regular Value}
A point \(y\in N\) is a \emph{regular value} if every point

\[
p\in f^{-1}(y)
\]

is a regular point.
\end{definitionbox}

If

\[
f^{-1}(y)=\varnothing,
\]

then \(y\) is automatically a regular value.

%%%%%%%%%%%%%%%%%%%%%%%%%%%%%%%%%%%%%%%%%%%%%%%%%%%%%%%%%%%%
\subsection{Regular Value Theorem}
\label{subsec:regular-value-theorem}
%%%%%%%%%%%%%%%%%%%%%%%%%%%%%%%%%%%%%%%%%%%%%%%%%%%%%%%%%%%%

One of the central results of differential topology is the Regular Value
Theorem.

\begin{theorem}[Regular Value Theorem]
Let

\[
f:M^m\to N^n
\]

be smooth, and let \(y\in N\) be a regular value.

Then

\[
\boxed{
f^{-1}(y)
}
\]

is a smooth submanifold of \(M\) of dimension

\[
\boxed{
m-n,
}
\]

provided the preimage is nonempty.
\end{theorem}

This theorem provides a powerful method for constructing manifolds as level
sets.

%%%%%%%%%%%%%%%%%%%%%%%%%%%%%%%%%%%%%%%%%%%%%%%%%%%%%%%%%%%%
\subsection{Worked Example: The Sphere as a Regular Level Set}
\label{subsec:sphere-regular-level-set}
%%%%%%%%%%%%%%%%%%%%%%%%%%%%%%%%%%%%%%%%%%%%%%%%%%%%%%%%%%%%

\begin{workedexamplebox}{Constructing the Sphere}

Define

\[
f:\R^{n+1}\to\R
\]

by

\[
f(x_1,\ldots,x_{n+1})
=
x_1^2+\cdots+x_{n+1}^2.
\]

Then

\[
S^n=f^{-1}(1).
\]

The gradient is

\[
\nabla f
=
(2x_1,\ldots,2x_{n+1}).
\]

On

\[
f^{-1}(1),
\]

the vector \(x\) cannot be zero.

Therefore

\[
\nabla f\neq0.
\]

Hence \(1\) is a regular value.

\begin{answerbox}
By the Regular Value Theorem,

\[
\boxed{
S^n
\text{ is a smooth }n\text{-dimensional submanifold of }\R^{n+1}.
}
\]
\end{answerbox}
\end{workedexamplebox}

%%%%%%%%%%%%%%%%%%%%%%%%%%%%%%%%%%%%%%%%%%%%%%%%%%%%%%%%%%%%
\subsection{Level Sets}
\label{subsec:differential-topology-level-sets}
%%%%%%%%%%%%%%%%%%%%%%%%%%%%%%%%%%%%%%%%%%%%%%%%%%%%%%%%%%%%

For a smooth function

\[
f:M\to\R,
\]

a set of the form

\[
\boxed{
f^{-1}(c)
}
\]

is called a level set.

When \(c\) is a regular value, the level set is a codimension-one
submanifold.

Thus if

\[
\dim M=n,
\]

then

\[
\boxed{
\dim f^{-1}(c)=n-1.
}
\]

%%%%%%%%%%%%%%%%%%%%%%%%%%%%%%%%%%%%%%%%%%%%%%%%%%%%%%%%%%%%
\subsection{Codimension}
\label{subsec:differential-topology-codimension}
%%%%%%%%%%%%%%%%%%%%%%%%%%%%%%%%%%%%%%%%%%%%%%%%%%%%%%%%%%%%

If

\[
S^k\subseteq M^n
\]

is a submanifold, its codimension is

\[
\boxed{
\operatorname{codim}(S)
=
n-k.
}
\]

For a regular level set of a real-valued function,

\[
\boxed{
\operatorname{codim}=1.
}
\]

%%%%%%%%%%%%%%%%%%%%%%%%%%%%%%%%%%%%%%%%%%%%%%%%%%%%%%%%%%%%
\subsection{Constant Rank Theorem}
\label{subsec:constant-rank-theorem}
%%%%%%%%%%%%%%%%%%%%%%%%%%%%%%%%%%%%%%%%%%%%%%%%%%%%%%%%%%%%

Suppose

\[
f:M^m\to N^n
\]

has constant rank \(r\) near a point.

Then suitable local coordinates can simplify \(f\) to the form

\[
\boxed{
(x_1,\ldots,x_m)
\mapsto
(x_1,\ldots,x_r,0,\ldots,0).
}
\]

This is the content of the Constant Rank Theorem.

It generalizes the local behavior described by the inverse and implicit
function theorems.

%%%%%%%%%%%%%%%%%%%%%%%%%%%%%%%%%%%%%%%%%%%%%%%%%%%%%%%%%%%%
\subsection{Inverse Function Theorem}
\label{subsec:differential-topology-inverse-function}
%%%%%%%%%%%%%%%%%%%%%%%%%%%%%%%%%%%%%%%%%%%%%%%%%%%%%%%%%%%%

Suppose

\[
f:M^n\to N^n
\]

and

\[
df_p
\]

is an isomorphism.

Then \(f\) is a local diffeomorphism near \(p\).

In Euclidean coordinates this corresponds to

\[
\boxed{
\det Df(p)\neq0.
}
\]

%%%%%%%%%%%%%%%%%%%%%%%%%%%%%%%%%%%%%%%%%%%%%%%%%%%%%%%%%%%%
\subsection{Implicit Function Theorem}
\label{subsec:differential-topology-implicit-function}
%%%%%%%%%%%%%%%%%%%%%%%%%%%%%%%%%%%%%%%%%%%%%%%%%%%%%%%%%%%%

The Implicit Function Theorem explains why regular level sets locally look
like Euclidean spaces.

For example, if

\[
F:\R^n\to\R^k
\]

has rank \(k\) at a point satisfying

\[
F(x)=c,
\]

then locally the solution set has dimension

\[
\boxed{
n-k.
}
\]

The Regular Value Theorem globalizes this local principle to manifolds.

%%%%%%%%%%%%%%%%%%%%%%%%%%%%%%%%%%%%%%%%%%%%%%%%%%%%%%%%%%%%
\subsection{Sard's Theorem}
\label{subsec:sards-theorem}
%%%%%%%%%%%%%%%%%%%%%%%%%%%%%%%%%%%%%%%%%%%%%%%%%%%%%%%%%%%%

Critical values might appear potentially complicated.

Sard's Theorem shows that they are nevertheless small in a measure-theoretic
sense.

\begin{theorem}[Sard's Theorem]
For a sufficiently differentiable map

\[
f:M\to N,
\]

the set of critical values of \(f\) has measure zero in \(N\).
\end{theorem}

In particular, for smooth maps,

\[
\boxed{
\text{almost every value is regular}
}
\]

in the measure-theoretic sense.

%%%%%%%%%%%%%%%%%%%%%%%%%%%%%%%%%%%%%%%%%%%%%%%%%%%%%%%%%%%%
\subsection{Why Sard's Theorem Matters}
\label{subsec:why-sard-matters}
%%%%%%%%%%%%%%%%%%%%%%%%%%%%%%%%%%%%%%%%%%%%%%%%%%%%%%%%%%%%

Sard's Theorem makes regular values abundant.

Therefore one can often choose a nearby value

\[
y
\]

such that

\[
f^{-1}(y)
\]

is a smooth submanifold.

This principle is fundamental in

\[
\boxed{
\text{transversality},
\quad
\text{intersection theory},
\quad
\text{degree theory},
\quad
\text{Morse theory}.
}
\]

%%%%%%%%%%%%%%%%%%%%%%%%%%%%%%%%%%%%%%%%%%%%%%%%%%%%%%%%%%%%
\subsection{Transverse Submanifolds}
\label{subsec:transverse-submanifolds}
%%%%%%%%%%%%%%%%%%%%%%%%%%%%%%%%%%%%%%%%%%%%%%%%%%%%%%%%%%%%

Let

\[
A,B\subseteq M
\]

be smooth submanifolds.

They intersect \emph{transversely} at

\[
p\in A\cap B
\]

if

\[
\boxed{
T_pA+T_pB=T_pM.
}
\]

We write

\[
\boxed{
A\pitchfork B.
}
\]

The symbol

\[
\pitchfork
\]

is read ``is transverse to.''

%%%%%%%%%%%%%%%%%%%%%%%%%%%%%%%%%%%%%%%%%%%%%%%%%%%%%%%%%%%%
\subsection{Dimension of a Transverse Intersection}
\label{subsec:dimension-transverse-intersection}
%%%%%%%%%%%%%%%%%%%%%%%%%%%%%%%%%%%%%%%%%%%%%%%%%%%%%%%%%%%%

If \(A\) and \(B\) intersect transversely, then

\[
A\cap B
\]

is a smooth submanifold.

Its dimension is

\[
\boxed{
\dim(A\cap B)
=
\dim A+\dim B-\dim M.
}
\]

Equivalently,

\[
\boxed{
\operatorname{codim}(A\cap B)
=
\operatorname{codim}(A)
+
\operatorname{codim}(B).
}
\]

%%%%%%%%%%%%%%%%%%%%%%%%%%%%%%%%%%%%%%%%%%%%%%%%%%%%%%%%%%%%
\subsection{Transverse Maps}
\label{subsec:transverse-maps}
%%%%%%%%%%%%%%%%%%%%%%%%%%%%%%%%%%%%%%%%%%%%%%%%%%%%%%%%%%%%

Let

\[
f:M\to N
\]

be smooth and let

\[
S\subseteq N
\]

be a submanifold.

We say \(f\) is transverse to \(S\), written

\[
\boxed{
f\pitchfork S,
}
\]

if for every

\[
p\in f^{-1}(S),
\]

we have

\[
\boxed{
df_p(T_pM)+T_{f(p)}S
=
T_{f(p)}N.
}
\]

%%%%%%%%%%%%%%%%%%%%%%%%%%%%%%%%%%%%%%%%%%%%%%%%%%%%%%%%%%%%
\subsection{Preimage Theorem}
\label{subsec:preimage-theorem}
%%%%%%%%%%%%%%%%%%%%%%%%%%%%%%%%%%%%%%%%%%%%%%%%%%%%%%%%%%%%

The Regular Value Theorem is a special case of a more general result.

\begin{theorem}[Preimage Theorem]
If

\[
f:M\to N
\]

is smooth and transverse to a submanifold

\[
S\subseteq N,
\]

then

\[
\boxed{
f^{-1}(S)
}
\]

is a smooth submanifold of \(M\).

Moreover,

\[
\boxed{
\operatorname{codim}f^{-1}(S)
=
\operatorname{codim}S.
}
\]
\end{theorem}

%%%%%%%%%%%%%%%%%%%%%%%%%%%%%%%%%%%%%%%%%%%%%%%%%%%%%%%%%%%%
\subsection{Genericity and Transversality}
\label{subsec:genericity-transversality}
%%%%%%%%%%%%%%%%%%%%%%%%%%%%%%%%%%%%%%%%%%%%%%%%%%%%%%%%%%%%

Transversality is important because transverse behavior is often generic.

A map that fails to be transverse can frequently be perturbed slightly into
one that is transverse.

Schematically,

\[
\boxed{
\text{arbitrary smooth map}
\quad
\overset{\text{small perturbation}}{\longrightarrow}
\quad
\text{transverse map}.
}
\]

This principle makes difficult intersections stable and manageable.

%%%%%%%%%%%%%%%%%%%%%%%%%%%%%%%%%%%%%%%%%%%%%%%%%%%%%%%%%%%%
\subsection{Transversality Theorem}
\label{subsec:transversality-theorem}
%%%%%%%%%%%%%%%%%%%%%%%%%%%%%%%%%%%%%%%%%%%%%%%%%%%%%%%%%%%%

One form of the Transversality Theorem states, roughly, that maps transverse
to a chosen submanifold form a large and generic collection among smooth
maps.

Precise formulations depend on the topology placed on the space of smooth
maps and on compactness assumptions.

The conceptual consequence is:

\[
\boxed{
\text{transverse intersections are the generic smooth situation}.
}
\]

%%%%%%%%%%%%%%%%%%%%%%%%%%%%%%%%%%%%%%%%%%%%%%%%%%%%%%%%%%%%
\subsection{Intersection Numbers}
\label{subsec:intersection-numbers}
%%%%%%%%%%%%%%%%%%%%%%%%%%%%%%%%%%%%%%%%%%%%%%%%%%%%%%%%%%%%

When oriented submanifolds intersect transversely in complementary
dimensions, their intersection consists locally of isolated points.

Each intersection point can be assigned a sign

\[
\boxed{
+1
\quad\text{or}\quad
-1.
}
\]

The signed sum gives an \emph{intersection number}.

This number often remains unchanged under suitable deformations.

%%%%%%%%%%%%%%%%%%%%%%%%%%%%%%%%%%%%%%%%%%%%%%%%%%%%%%%%%%%%
\subsection{Orientation}
\label{subsec:differential-topology-orientation}
%%%%%%%%%%%%%%%%%%%%%%%%%%%%%%%%%%%%%%%%%%%%%%%%%%%%%%%%%%%%

An orientation of an \(n\)-dimensional manifold provides a consistent choice
of orientation for its tangent spaces.

In coordinates, orientation is related to ordered bases.

Two ordered bases have the same orientation when the change-of-basis matrix
has positive determinant.

Thus

\[
\boxed{
\det A>0
}
\]

preserves orientation, while

\[
\boxed{
\det A<0
}
\]

reverses orientation.

%%%%%%%%%%%%%%%%%%%%%%%%%%%%%%%%%%%%%%%%%%%%%%%%%%%%%%%%%%%%
\subsection{Orientable Manifolds}
\label{subsec:orientable-manifolds}
%%%%%%%%%%%%%%%%%%%%%%%%%%%%%%%%%%%%%%%%%%%%%%%%%%%%%%%%%%%%

\begin{definitionbox}{Orientable Manifold}
A smooth manifold is \emph{orientable} if a consistent orientation can be
chosen throughout the manifold.
\end{definitionbox}

Examples include

\[
\boxed{
S^n,
\qquad
T^n,
\qquad
\R^n.
}
\]

The M\"obius strip is a standard example of a nonorientable surface with
boundary.

The real projective plane

\[
\mathbb{RP}^2
\]

is a closed nonorientable surface.

%%%%%%%%%%%%%%%%%%%%%%%%%%%%%%%%%%%%%%%%%%%%%%%%%%%%%%%%%%%%
\subsection{Manifolds with Boundary}
\label{subsec:manifolds-with-boundary-differential}
%%%%%%%%%%%%%%%%%%%%%%%%%%%%%%%%%%%%%%%%%%%%%%%%%%%%%%%%%%%%

A manifold with boundary locally resembles either

\[
\R^n
\]

or the half-space

\[
\boxed{
\mathbb{H}^n
=
\{
(x_1,\ldots,x_n)\in\R^n:x_n\geq0
\}.
}
\]

Points modeled on

\[
x_n=0
\]

form the boundary

\[
\boxed{
\partial M.
}
\]

%%%%%%%%%%%%%%%%%%%%%%%%%%%%%%%%%%%%%%%%%%%%%%%%%%%%%%%%%%%%
\subsection{Boundary of a Manifold}
\label{subsec:boundary-manifold-differential}
%%%%%%%%%%%%%%%%%%%%%%%%%%%%%%%%%%%%%%%%%%%%%%%%%%%%%%%%%%%%

If \(M\) is an \(n\)-dimensional smooth manifold with boundary, then

\[
\boxed{
\partial M
}
\]

is an \((n-1)\)-dimensional smooth manifold without boundary.

Furthermore,

\[
\boxed{
\partial(\partial M)=\varnothing.
}
\]

This is the smooth-manifold version of the principle

\[
\boxed{
\text{the boundary of a boundary is zero}.
}
\]

%%%%%%%%%%%%%%%%%%%%%%%%%%%%%%%%%%%%%%%%%%%%%%%%%%%%%%%%%%%%
\subsection{Boundary Orientation}
\label{subsec:boundary-orientation}
%%%%%%%%%%%%%%%%%%%%%%%%%%%%%%%%%%%%%%%%%%%%%%%%%%%%%%%%%%%%

If \(M\) is oriented, its boundary inherits a natural orientation.

A common convention is the ``outward normal first'' rule.

If

\[
\nu
\]

is an outward-pointing normal vector, then

\[
(\nu,v_1,\ldots,v_{n-1})
\]

is positively oriented in \(M\) precisely when

\[
(v_1,\ldots,v_{n-1})
\]

is positively oriented in

\[
\partial M.
\]

%%%%%%%%%%%%%%%%%%%%%%%%%%%%%%%%%%%%%%%%%%%%%%%%%%%%%%%%%%%%
\subsection{Degree of a Smooth Map}
\label{subsec:differential-topology-degree}
%%%%%%%%%%%%%%%%%%%%%%%%%%%%%%%%%%%%%%%%%%%%%%%%%%%%%%%%%%%%

Let

\[
f:M^n\to N^n
\]

be a smooth map between compact connected oriented \(n\)-manifolds without
boundary.

Choose a regular value

\[
y\in N.
\]

For each

\[
x\in f^{-1}(y),
\]

the derivative

\[
df_x
\]

is an isomorphism.

Assign

\[
\operatorname{sgn}(df_x)
=
\begin{cases}
+1,
&
df_x\text{ preserves orientation},
\\
-1,
&
df_x\text{ reverses orientation}.
\end{cases}
\]

Then

\[
\boxed{
\deg(f)
=
\sum_{x\in f^{-1}(y)}
\operatorname{sgn}(df_x).
}
\]

%%%%%%%%%%%%%%%%%%%%%%%%%%%%%%%%%%%%%%%%%%%%%%%%%%%%%%%%%%%%
\subsection{Properties of Degree}
\label{subsec:degree-properties-differential}
%%%%%%%%%%%%%%%%%%%%%%%%%%%%%%%%%%%%%%%%%%%%%%%%%%%%%%%%%%%%

Degree is independent of the chosen regular value under the standard
hypotheses.

It is also homotopy invariant:

\[
\boxed{
f\simeq g
\Longrightarrow
\deg(f)=\deg(g).
}
\]

Furthermore,

\[
\boxed{
\deg(g\circ f)
=
\deg(g)\deg(f).
}
\]

If

\[
\deg(f)\neq0,
\]

then \(f\) must be surjective.

%%%%%%%%%%%%%%%%%%%%%%%%%%%%%%%%%%%%%%%%%%%%%%%%%%%%%%%%%%%%
\subsection{Mod-\texorpdfstring{\(2\)}{2} Degree}
\label{subsec:mod-two-degree}
%%%%%%%%%%%%%%%%%%%%%%%%%%%%%%%%%%%%%%%%%%%%%%%%%%%%%%%%%%%%

When orientations are unavailable, one can count preimages modulo \(2\).

For a regular value \(y\),

\[
\boxed{
\deg_2(f)
=
\#f^{-1}(y)\pmod 2.
}
\]

The notation

\[
\deg_2
\]

is read ``degree mod two.''

This provides a useful invariant even for nonorientable situations.

%%%%%%%%%%%%%%%%%%%%%%%%%%%%%%%%%%%%%%%%%%%%%%%%%%%%%%%%%%%%
\subsection{Vector Fields}
\label{subsec:differential-topology-vector-fields}
%%%%%%%%%%%%%%%%%%%%%%%%%%%%%%%%%%%%%%%%%%%%%%%%%%%%%%%%%%%%

A smooth vector field assigns a tangent vector to every point:

\[
\boxed{
X(p)\in T_pM.
}
\]

Formally, it is a smooth section

\[
\boxed{
X:M\to TM
}
\]

of the tangent bundle satisfying

\[
\pi\circ X=\operatorname{id}_M,
\]

where

\[
\pi:TM\to M
\]

is the tangent-bundle projection.

%%%%%%%%%%%%%%%%%%%%%%%%%%%%%%%%%%%%%%%%%%%%%%%%%%%%%%%%%%%%
\subsection{Zeros of Vector Fields}
\label{subsec:zeros-vector-fields}
%%%%%%%%%%%%%%%%%%%%%%%%%%%%%%%%%%%%%%%%%%%%%%%%%%%%%%%%%%%%

A point

\[
p\in M
\]

is a zero of a vector field \(X\) if

\[
\boxed{
X(p)=0.
}
\]

Zeros of vector fields encode important topological information.

Locally they can be assigned indices.

Globally the sum of these indices is related to the Euler characteristic.

%%%%%%%%%%%%%%%%%%%%%%%%%%%%%%%%%%%%%%%%%%%%%%%%%%%%%%%%%%%%
\subsection{Hairy Ball Theorem}
\label{subsec:hairy-ball-theorem}
%%%%%%%%%%%%%%%%%%%%%%%%%%%%%%%%%%%%%%%%%%%%%%%%%%%%%%%%%%%%

\begin{theorem}[Hairy Ball Theorem]
Every continuous tangent vector field on an even-dimensional sphere

\[
S^{2n}
\]

must vanish somewhere.
\end{theorem}

For \(S^2\), this means there is no continuous nowhere-zero tangent vector
field.

The common phrase is:

\[
\boxed{
\text{you cannot comb a hairy sphere flat without creating a cowlick}.
}
\]

%%%%%%%%%%%%%%%%%%%%%%%%%%%%%%%%%%%%%%%%%%%%%%%%%%%%%%%%%%%%
\subsection{Odd-Dimensional Spheres}
\label{subsec:odd-dimensional-sphere-vector-fields}
%%%%%%%%%%%%%%%%%%%%%%%%%%%%%%%%%%%%%%%%%%%%%%%%%%%%%%%%%%%%

Odd-dimensional spheres behave differently.

For example, on

\[
S^1
\]

there is an obvious nowhere-zero tangent vector field.

Similarly,

\[
S^3
\]

admits nowhere-zero tangent vector fields.

Thus the obstruction in the Hairy Ball Theorem depends fundamentally on
topology and dimension.

%%%%%%%%%%%%%%%%%%%%%%%%%%%%%%%%%%%%%%%%%%%%%%%%%%%%%%%%%%%%
\subsection{Index of a Vector Field}
\label{subsec:index-vector-field}
%%%%%%%%%%%%%%%%%%%%%%%%%%%%%%%%%%%%%%%%%%%%%%%%%%%%%%%%%%%%

Suppose \(p\) is an isolated zero of a vector field.

One can study the direction of the vector field on a small sphere
surrounding \(p\).

After normalization, this produces a map

\[
S^{n-1}\to S^{n-1}.
\]

The degree of this map is called the \emph{index} of the vector field at
\(p\):

\[
\boxed{
\operatorname{ind}_p(X).
}
\]

%%%%%%%%%%%%%%%%%%%%%%%%%%%%%%%%%%%%%%%%%%%%%%%%%%%%%%%%%%%%
\subsection{Poincar\'e--Hopf Theorem}
\label{subsec:poincare-hopf}
%%%%%%%%%%%%%%%%%%%%%%%%%%%%%%%%%%%%%%%%%%%%%%%%%%%%%%%%%%%%

\begin{theorem}[Poincar\'e--Hopf Theorem]
Let \(M\) be a compact smooth manifold without boundary, and let \(X\) be a
smooth vector field with isolated zeros.

Then

\[
\boxed{
\sum_{p:X(p)=0}
\operatorname{ind}_p(X)
=
\chi(M).
}
\]
\end{theorem}

This theorem creates a remarkable bridge:

\[
\boxed{
\text{local behavior of a vector field}
\longleftrightarrow
\text{global topology}.
}
\]

%%%%%%%%%%%%%%%%%%%%%%%%%%%%%%%%%%%%%%%%%%%%%%%%%%%%%%%%%%%%
\subsection{Example: Vector Fields on the Torus}
\label{subsec:vector-fields-torus}
%%%%%%%%%%%%%%%%%%%%%%%%%%%%%%%%%%%%%%%%%%%%%%%%%%%%%%%%%%%%

Since

\[
\chi(T^2)=0,
\]

the torus admits nowhere-zero tangent vector fields.

This contrasts with

\[
S^2,
\]

for which

\[
\chi(S^2)=2.
\]

The Euler characteristic therefore provides an obstruction to nowhere-zero
vector fields.

%%%%%%%%%%%%%%%%%%%%%%%%%%%%%%%%%%%%%%%%%%%%%%%%%%%%%%%%%%%%
\subsection{Homotopy of Smooth Maps}
\label{subsec:smooth-homotopy}
%%%%%%%%%%%%%%%%%%%%%%%%%%%%%%%%%%%%%%%%%%%%%%%%%%%%%%%%%%%%

A smooth homotopy between smooth maps

\[
f,g:M\to N
\]

is a smooth map

\[
\boxed{
H:M\times[0,1]\to N
}
\]

with

\[
H(x,0)=f(x)
\]

and

\[
H(x,1)=g(x).
\]

Smooth homotopy is analogous to ordinary homotopy but respects the
differentiable structure.

%%%%%%%%%%%%%%%%%%%%%%%%%%%%%%%%%%%%%%%%%%%%%%%%%%%%%%%%%%%%
\subsection{Isotopy}
\label{subsec:isotopy}
%%%%%%%%%%%%%%%%%%%%%%%%%%%%%%%%%%%%%%%%%%%%%%%%%%%%%%%%%%%%

\begin{definitionbox}{Isotopy}
An \emph{isotopy} is a smooth one-parameter family of embeddings

\[
\boxed{
F_t:M\to N,
\qquad
0\leq t\leq1.
}
\]
\end{definitionbox}

Thus every intermediate map

\[
F_t
\]

must remain an embedding.

This is much stronger than an arbitrary homotopy.

%%%%%%%%%%%%%%%%%%%%%%%%%%%%%%%%%%%%%%%%%%%%%%%%%%%%%%%%%%%%
\subsection{Isotopy Versus Homotopy}
\label{subsec:isotopy-vs-homotopy}
%%%%%%%%%%%%%%%%%%%%%%%%%%%%%%%%%%%%%%%%%%%%%%%%%%%%%%%%%%%%

A homotopy may pass through maps with

\[
\boxed{
\text{self-intersections},
\quad
\text{singularities},
\quad
\text{collapsed regions}.
}
\]

An isotopy cannot.

Therefore isotopy captures the idea of smoothly deforming an embedded object
without tearing it or allowing it to pass through itself.

This concept is fundamental in knot theory.

%%%%%%%%%%%%%%%%%%%%%%%%%%%%%%%%%%%%%%%%%%%%%%%%%%%%%%%%%%%%
\subsection{Ambient Isotopy}
\label{subsec:ambient-isotopy}
%%%%%%%%%%%%%%%%%%%%%%%%%%%%%%%%%%%%%%%%%%%%%%%%%%%%%%%%%%%%

An \emph{ambient isotopy} deforms the surrounding space itself.

It is a smooth family of diffeomorphisms

\[
\boxed{
F_t:N\to N
}
\]

with

\[
F_0=\operatorname{id}_N.
\]

Two embedded submanifolds are ambient isotopic if one can be carried to the
other by such a deformation.

For knots, ambient isotopy gives the standard notion of knot equivalence.

%%%%%%%%%%%%%%%%%%%%%%%%%%%%%%%%%%%%%%%%%%%%%%%%%%%%%%%%%%%%
\subsection{Isotopy Extension Principle}
\label{subsec:isotopy-extension}
%%%%%%%%%%%%%%%%%%%%%%%%%%%%%%%%%%%%%%%%%%%%%%%%%%%%%%%%%%%%

Under standard compactness and manifold hypotheses, an isotopy of an
embedded submanifold can often be extended to an ambient isotopy of the
surrounding manifold.

This is known as an \emph{Isotopy Extension Theorem}.

It provides an important bridge between deformation of a submanifold and
deformation of the ambient space.

%%%%%%%%%%%%%%%%%%%%%%%%%%%%%%%%%%%%%%%%%%%%%%%%%%%%%%%%%%%%
\subsection{Morse Theory}
\label{subsec:morse-theory-introduction}
%%%%%%%%%%%%%%%%%%%%%%%%%%%%%%%%%%%%%%%%%%%%%%%%%%%%%%%%%%%%

Morse theory studies the topology of a smooth manifold using critical points
of smooth real-valued functions.

Consider

\[
\boxed{
f:M\to\R.
}
\]

The topology of the level sets and sublevel sets of \(f\) changes when the
value passes through critical values.

This converts calculus into topology.

%%%%%%%%%%%%%%%%%%%%%%%%%%%%%%%%%%%%%%%%%%%%%%%%%%%%%%%%%%%%
\subsection{Critical Points of Real-Valued Functions}
\label{subsec:morse-critical-points}
%%%%%%%%%%%%%%%%%%%%%%%%%%%%%%%%%%%%%%%%%%%%%%%%%%%%%%%%%%%%

A point

\[
p\in M
\]

is critical when

\[
\boxed{
df_p=0.
}
\]

In local coordinates,

\[
\boxed{
\nabla f(p)=0.
}
\]

Critical points include familiar objects from multivariable calculus such as

\[
\boxed{
\text{local minima},
\quad
\text{local maxima},
\quad
\text{saddle points}.
}
\]

%%%%%%%%%%%%%%%%%%%%%%%%%%%%%%%%%%%%%%%%%%%%%%%%%%%%%%%%%%%%
\subsection{Hessian}
\label{subsec:morse-hessian}
%%%%%%%%%%%%%%%%%%%%%%%%%%%%%%%%%%%%%%%%%%%%%%%%%%%%%%%%%%%%

At a critical point, the second-order behavior is encoded by the Hessian.

In local coordinates,

\[
\boxed{
\operatorname{Hess}_p(f)
=
\left[
\frac{\partial^2f}
{\partial x_i\partial x_j}(p)
\right].
}
\]

At a critical point, the associated quadratic form has a coordinate-independent
meaning.

%%%%%%%%%%%%%%%%%%%%%%%%%%%%%%%%%%%%%%%%%%%%%%%%%%%%%%%%%%%%
\subsection{Nondegenerate Critical Points}
\label{subsec:nondegenerate-critical-points}
%%%%%%%%%%%%%%%%%%%%%%%%%%%%%%%%%%%%%%%%%%%%%%%%%%%%%%%%%%%%

\begin{definitionbox}{Nondegenerate Critical Point}
A critical point \(p\) is \emph{nondegenerate} if its Hessian is
nonsingular:

\[
\boxed{
\det\operatorname{Hess}_p(f)\neq0.
}
\]
\end{definitionbox}

Nondegenerate critical points are stable under sufficiently small
perturbations.

%%%%%%%%%%%%%%%%%%%%%%%%%%%%%%%%%%%%%%%%%%%%%%%%%%%%%%%%%%%%
\subsection{Morse Functions}
\label{subsec:morse-functions}
%%%%%%%%%%%%%%%%%%%%%%%%%%%%%%%%%%%%%%%%%%%%%%%%%%%%%%%%%%%%

\begin{definitionbox}{Morse Function}
A smooth function

\[
f:M\to\R
\]

is a \emph{Morse function} if every critical point of \(f\) is
nondegenerate.
\end{definitionbox}

Morse functions are generic among smooth real-valued functions in an
appropriate sense.

%%%%%%%%%%%%%%%%%%%%%%%%%%%%%%%%%%%%%%%%%%%%%%%%%%%%%%%%%%%%
\subsection{Morse Index}
\label{subsec:morse-index}
%%%%%%%%%%%%%%%%%%%%%%%%%%%%%%%%%%%%%%%%%%%%%%%%%%%%%%%%%%%%

At a nondegenerate critical point, the Hessian has no zero eigenvalues.

\begin{definitionbox}{Morse Index}
The \emph{Morse index} of a nondegenerate critical point \(p\) is the number
of negative eigenvalues of

\[
\operatorname{Hess}_p(f),
\]

counted with multiplicity.
\end{definitionbox}

It is usually denoted

\[
\boxed{
\lambda(p).
}
\]

The lowercase Greek letter

\[
\lambda
\]

is read ``lambda.''

%%%%%%%%%%%%%%%%%%%%%%%%%%%%%%%%%%%%%%%%%%%%%%%%%%%%%%%%%%%%
\subsection{Meaning of Morse Index}
\label{subsec:meaning-morse-index}
%%%%%%%%%%%%%%%%%%%%%%%%%%%%%%%%%%%%%%%%%%%%%%%%%%%%%%%%%%%%

For a function on an \(n\)-manifold:

\[
\lambda=0
\]

corresponds locally to a minimum.

\[
\lambda=n
\]

corresponds locally to a maximum.

Intermediate values correspond to saddle-type critical points.

Thus the index measures the number of independent downward directions.

%%%%%%%%%%%%%%%%%%%%%%%%%%%%%%%%%%%%%%%%%%%%%%%%%%%%%%%%%%%%
\subsection{Morse Lemma}
\label{subsec:morse-lemma}
%%%%%%%%%%%%%%%%%%%%%%%%%%%%%%%%%%%%%%%%%%%%%%%%%%%%%%%%%%%%

\begin{theorem}[Morse Lemma]
Let \(p\) be a nondegenerate critical point of a smooth function

\[
f:M^n\to\R
\]

with Morse index \(\lambda\).

Then there exist local coordinates centered at \(p\) such that

\[
\boxed{
f(x)
=
f(p)
-
x_1^2-\cdots-x_\lambda^2
+
x_{\lambda+1}^2+\cdots+x_n^2.
}
\]
\end{theorem}

Thus every nondegenerate critical point has a standard local form.

%%%%%%%%%%%%%%%%%%%%%%%%%%%%%%%%%%%%%%%%%%%%%%%%%%%%%%%%%%%%
\subsection{Sublevel Sets}
\label{subsec:morse-sublevel-sets}
%%%%%%%%%%%%%%%%%%%%%%%%%%%%%%%%%%%%%%%%%%%%%%%%%%%%%%%%%%%%

For

\[
a\in\R,
\]

define the sublevel set

\[
\boxed{
M_a
=
f^{-1}((-\infty,a]).
}
\]

Morse theory studies how

\[
M_a
\]

changes as \(a\) increases.

If the interval

\[
[a,b]
\]

contains no critical values, then under suitable compactness hypotheses,

\[
\boxed{
M_a
\text{ and }
M_b
}
\]

have the same relevant smooth/topological type.

%%%%%%%%%%%%%%%%%%%%%%%%%%%%%%%%%%%%%%%%%%%%%%%%%%%%%%%%%%%%
\subsection{Passing a Critical Value}
\label{subsec:passing-critical-value}
%%%%%%%%%%%%%%%%%%%%%%%%%%%%%%%%%%%%%%%%%%%%%%%%%%%%%%%%%%%%

Suppose \(c\) is a critical value corresponding to one nondegenerate
critical point of index \(\lambda\).

As the sublevel set passes through \(c\), its topology changes by attaching
a \(\lambda\)-dimensional handle.

Schematically,

\[
\boxed{
M_{c-\varepsilon}
\longrightarrow
M_{c+\varepsilon}
}
\]

is accomplished by attaching a handle of index \(\lambda\).

%%%%%%%%%%%%%%%%%%%%%%%%%%%%%%%%%%%%%%%%%%%%%%%%%%%%%%%%%%%%
\subsection{Handles}
\label{subsec:handles-differential-topology}
%%%%%%%%%%%%%%%%%%%%%%%%%%%%%%%%%%%%%%%%%%%%%%%%%%%%%%%%%%%%

An \(n\)-dimensional handle of index \(\lambda\) is

\[
\boxed{
D^\lambda\times D^{n-\lambda}.
}
\]

It is attached along part of its boundary,

\[
\boxed{
S^{\lambda-1}\times D^{n-\lambda}.
}
\]

Handles provide a geometric way to build manifolds from simpler pieces.

%%%%%%%%%%%%%%%%%%%%%%%%%%%%%%%%%%%%%%%%%%%%%%%%%%%%%%%%%%%%
\subsection{Morse Theory on the Sphere}
\label{subsec:morse-sphere}
%%%%%%%%%%%%%%%%%%%%%%%%%%%%%%%%%%%%%%%%%%%%%%%%%%%%%%%%%%%%

Consider the height function

\[
f:S^2\to\R,
\qquad
f(x,y,z)=z.
\]

There are two critical points:

\[
\boxed{
\text{south pole}
}
\]

and

\[
\boxed{
\text{north pole}.
}
\]

The south pole has index \(0\).

The north pole has index \(2\).

The sphere can therefore be understood using one \(0\)-handle and one
\(2\)-handle.

%%%%%%%%%%%%%%%%%%%%%%%%%%%%%%%%%%%%%%%%%%%%%%%%%%%%%%%%%%%%
\subsection{Morse Theory on the Torus}
\label{subsec:morse-torus}
%%%%%%%%%%%%%%%%%%%%%%%%%%%%%%%%%%%%%%%%%%%%%%%%%%%%%%%%%%%%

A generic height function on a vertically positioned torus has four critical
points:

\[
\boxed{
\text{one minimum},
}
\]

\[
\boxed{
\text{two saddle points},
}
\]

\[
\boxed{
\text{one maximum}.
}
\]

Their Morse indices are

\[
\boxed{
0,\quad1,\quad1,\quad2.
}
\]

These critical points reflect the topology of the torus.

%%%%%%%%%%%%%%%%%%%%%%%%%%%%%%%%%%%%%%%%%%%%%%%%%%%%%%%%%%%%
\subsection{Morse Numbers and Euler Characteristic}
\label{subsec:morse-euler}
%%%%%%%%%%%%%%%%%%%%%%%%%%%%%%%%%%%%%%%%%%%%%%%%%%%%%%%%%%%%

Let

\[
m_k
\]

denote the number of critical points of Morse index \(k\).

Then

\[
\boxed{
\chi(M)
=
\sum_k(-1)^k m_k.
}
\]

For the torus,

\[
m_0=1,
\qquad
m_1=2,
\qquad
m_2=1.
\]

Therefore

\[
\boxed{
\chi(T^2)
=
1-2+1
=
0.
}
\]

%%%%%%%%%%%%%%%%%%%%%%%%%%%%%%%%%%%%%%%%%%%%%%%%%%%%%%%%%%%%
\subsection{Morse Inequalities}
\label{subsec:morse-inequalities}
%%%%%%%%%%%%%%%%%%%%%%%%%%%%%%%%%%%%%%%%%%%%%%%%%%%%%%%%%%%%

Morse theory relates numbers of critical points to Betti numbers.

In particular,

\[
\boxed{
m_k\geq b_k
}
\]

is part of the weak Morse inequalities.

Thus topology forces smooth functions to possess certain critical points.

A manifold with complicated homology cannot admit a Morse function with
arbitrarily few critical points.

%%%%%%%%%%%%%%%%%%%%%%%%%%%%%%%%%%%%%%%%%%%%%%%%%%%%%%%%%%%%
\subsection{Morse Theory and Homology}
\label{subsec:morse-homology-connection}
%%%%%%%%%%%%%%%%%%%%%%%%%%%%%%%%%%%%%%%%%%%%%%%%%%%%%%%%%%%%

Morse theory can be developed further into \emph{Morse homology}.

One constructs a chain complex generated by critical points of a Morse
function.

The boundary operator counts suitable gradient-flow trajectories between
critical points.

The resulting homology is isomorphic to the ordinary homology of the
manifold.

Thus

\[
\boxed{
\text{critical points}
+
\text{dynamics}
\longrightarrow
\text{topological homology}.
}
\]

%%%%%%%%%%%%%%%%%%%%%%%%%%%%%%%%%%%%%%%%%%%%%%%%%%%%%%%%%%%%
\subsection{Cobordism}
\label{subsec:cobordism-introduction}
%%%%%%%%%%%%%%%%%%%%%%%%%%%%%%%%%%%%%%%%%%%%%%%%%%%%%%%%%%%%

Differential topology also studies when manifolds can occur as boundaries of
higher-dimensional manifolds.

\begin{definitionbox}{Cobordism}
Two closed \(n\)-manifolds \(M_0\) and \(M_1\) are \emph{cobordant} if there
exists a compact \((n+1)\)-manifold \(W\) such that

\[
\boxed{
\partial W
\cong
M_0\sqcup M_1
}
\]

in the unoriented setting, with the appropriate orientation reversal in the
oriented setting.
\end{definitionbox}

The symbol

\[
\sqcup
\]

is read ``disjoint union.''

%%%%%%%%%%%%%%%%%%%%%%%%%%%%%%%%%%%%%%%%%%%%%%%%%%%%%%%%%%%%
\subsection{Oriented Cobordism}
\label{subsec:oriented-cobordism}
%%%%%%%%%%%%%%%%%%%%%%%%%%%%%%%%%%%%%%%%%%%%%%%%%%%%%%%%%%%%

For oriented manifolds, the boundary orientation matters.

One writes

\[
\boxed{
\partial W
=
M_1\sqcup(-M_0),
}
\]

where

\[
-M_0
\]

means \(M_0\) with reversed orientation.

Cobordism becomes an equivalence relation and leads to cobordism groups.

%%%%%%%%%%%%%%%%%%%%%%%%%%%%%%%%%%%%%%%%%%%%%%%%%%%%%%%%%%%%
\subsection{Example of a Cobordism}
\label{subsec:cylinder-cobordism}
%%%%%%%%%%%%%%%%%%%%%%%%%%%%%%%%%%%%%%%%%%%%%%%%%%%%%%%%%%%%

The cylinder

\[
\boxed{
S^1\times[0,1]
}
\]

is a cobordism between two circles.

Its boundary is

\[
\boxed{
\partial(S^1\times[0,1])
=
(S^1\times\{1\})
\sqcup
-(S^1\times\{0\})
}
\]

with the induced orientations.

%%%%%%%%%%%%%%%%%%%%%%%%%%%%%%%%%%%%%%%%%%%%%%%%%%%%%%%%%%%%
\subsection{Cobordism and Morse Theory}
\label{subsec:cobordism-morse-theory}
%%%%%%%%%%%%%%%%%%%%%%%%%%%%%%%%%%%%%%%%%%%%%%%%%%%%%%%%%%%%

Morse functions on cobordisms reveal how one manifold can transform into
another through handle attachments.

Thus Morse theory provides a geometric language for decomposing cobordisms.

This connection is fundamental in high-dimensional manifold topology.

%%%%%%%%%%%%%%%%%%%%%%%%%%%%%%%%%%%%%%%%%%%%%%%%%%%%%%%%%%%%
\subsection{Smooth Structures}
\label{subsec:smooth-structures}
%%%%%%%%%%%%%%%%%%%%%%%%%%%%%%%%%%%%%%%%%%%%%%%%%%%%%%%%%%%%

A topological manifold may admit more than one inequivalent smooth
structure.

Two smooth structures are considered equivalent if there is a
diffeomorphism between them.

In many dimensions the relationship between topological and smooth
classification is highly nontrivial.

%%%%%%%%%%%%%%%%%%%%%%%%%%%%%%%%%%%%%%%%%%%%%%%%%%%%%%%%%%%%
\subsection{Exotic Spheres}
\label{subsec:exotic-spheres}
%%%%%%%%%%%%%%%%%%%%%%%%%%%%%%%%%%%%%%%%%%%%%%%%%%%%%%%%%%%%

An \emph{exotic sphere} is a smooth manifold that is homeomorphic to a
standard sphere

\[
S^n
\]

but not diffeomorphic to it.

Thus an exotic sphere has the same underlying topological type as \(S^n\)
but a different smooth structure.

Exotic spheres demonstrate that

\[
\boxed{
\text{topological equivalence}
\not\Rightarrow
\text{smooth equivalence}.
}
\]

%%%%%%%%%%%%%%%%%%%%%%%%%%%%%%%%%%%%%%%%%%%%%%%%%%%%%%%%%%%%
\subsection{Exotic Euclidean Space}
\label{subsec:exotic-r4}
%%%%%%%%%%%%%%%%%%%%%%%%%%%%%%%%%%%%%%%%%%%%%%%%%%%%%%%%%%%%

A particularly remarkable phenomenon occurs in dimension four.

There exist smooth manifolds homeomorphic to

\[
\R^4
\]

that are not diffeomorphic to the standard

\[
\R^4.
\]

These are called \emph{exotic \(\R^4\)'s}.

Dimension four is exceptional in the theory of smooth manifolds.

%%%%%%%%%%%%%%%%%%%%%%%%%%%%%%%%%%%%%%%%%%%%%%%%%%%%%%%%%%%%
\subsection{Differential Topology and Knot Theory}
\label{subsec:differential-topology-knot-theory}
%%%%%%%%%%%%%%%%%%%%%%%%%%%%%%%%%%%%%%%%%%%%%%%%%%%%%%%%%%%%

A smooth knot may be viewed as a smooth embedding

\[
\boxed{
S^1\hookrightarrow\R^3
}
\]

or

\[
S^1\hookrightarrow S^3.
\]

Two knots are considered equivalent when their embeddings are related by
ambient isotopy.

Thus knot theory naturally belongs to both differential and geometric
topology.

%%%%%%%%%%%%%%%%%%%%%%%%%%%%%%%%%%%%%%%%%%%%%%%%%%%%%%%%%%%%
\subsection{Differential Topology and Dynamical Systems}
\label{subsec:differential-topology-dynamics}
%%%%%%%%%%%%%%%%%%%%%%%%%%%%%%%%%%%%%%%%%%%%%%%%%%%%%%%%%%%%

Vector fields generate differential equations and flows.

Zeros of vector fields correspond to equilibrium points.

Differential topology constrains which vector fields can exist globally.

For example, the Hairy Ball Theorem shows that certain manifolds cannot
support nowhere-zero tangent vector fields.

Thus

\[
\boxed{
\text{topology constrains dynamics}.
}
\]

%%%%%%%%%%%%%%%%%%%%%%%%%%%%%%%%%%%%%%%%%%%%%%%%%%%%%%%%%%%%
\subsection{Differential Topology and Optimization}
\label{subsec:differential-topology-optimization}
%%%%%%%%%%%%%%%%%%%%%%%%%%%%%%%%%%%%%%%%%%%%%%%%%%%%%%%%%%%%

Optimization studies critical points of functions.

Differential topology places these critical points into a global framework.

Morse theory shows that the number and type of critical points are related
to the topology of the underlying manifold.

Therefore the geometry of an optimization landscape is not independent of
the topology of its domain.

%%%%%%%%%%%%%%%%%%%%%%%%%%%%%%%%%%%%%%%%%%%%%%%%%%%%%%%%%%%%
\subsection{Differential Topology and Physics}
\label{subsec:differential-topology-physics}
%%%%%%%%%%%%%%%%%%%%%%%%%%%%%%%%%%%%%%%%%%%%%%%%%%%%%%%%%%%%

Smooth manifolds, vector fields, degree theory, and characteristic
structures appear throughout mathematical physics.

Applications include

\[
\boxed{
\text{classical mechanics},
\quad
\text{relativity},
\quad
\text{gauge theory},
\quad
\text{field theory},
\quad
\text{topological defects}.
}
\]

Many physical fields are mathematically sections of bundles over smooth
manifolds.

%%%%%%%%%%%%%%%%%%%%%%%%%%%%%%%%%%%%%%%%%%%%%%%%%%%%%%%%%%%%
\subsection{Worked Example: Regular Value}
\label{subsec:worked-regular-value}
%%%%%%%%%%%%%%%%%%%%%%%%%%%%%%%%%%%%%%%%%%%%%%%%%%%%%%%%%%%%

\begin{workedexamplebox}{A Hyperbola as a Smooth Submanifold}

Consider

\[
f:\R^2\to\R,
\qquad
f(x,y)=xy.
\]

Determine whether

\[
f^{-1}(1)
\]

is a smooth submanifold.

Compute

\[
\nabla f(x,y)
=
(y,x).
\]

On

\[
f^{-1}(1),
\]

we have

\[
xy=1.
\]

Therefore neither \(x\) nor \(y\) is zero.

Hence

\[
\nabla f(x,y)\neq(0,0).
\]

Thus \(1\) is a regular value.

\begin{answerbox}
\[
\boxed{
f^{-1}(1)
=
\{(x,y):xy=1\}
}
\]

is a smooth \(1\)-dimensional submanifold of \(\R^2\).
\end{answerbox}
\end{workedexamplebox}

%%%%%%%%%%%%%%%%%%%%%%%%%%%%%%%%%%%%%%%%%%%%%%%%%%%%%%%%%%%%
\subsection{Worked Example: Critical Points and Morse Index}
\label{subsec:worked-morse-index}
%%%%%%%%%%%%%%%%%%%%%%%%%%%%%%%%%%%%%%%%%%%%%%%%%%%%%%%%%%%%

\begin{workedexamplebox}{Classifying a Critical Point}

Consider

\[
f(x,y)
=
x^2-y^2.
\]

The gradient is

\[
\nabla f
=
(2x,-2y).
\]

Thus the only critical point is

\[
(0,0).
\]

The Hessian is

\[
\operatorname{Hess}(f)
=
\begin{pmatrix}
2&0\\
0&-2
\end{pmatrix}.
\]

Its eigenvalues are

\[
2
\qquad\text{and}\qquad
-2.
\]

There is exactly one negative eigenvalue.

\begin{answerbox}
\[
\boxed{
(0,0)
\text{ is a nondegenerate critical point of Morse index }1.
}
\]

It is a saddle point.
\end{answerbox}
\end{workedexamplebox}

%%%%%%%%%%%%%%%%%%%%%%%%%%%%%%%%%%%%%%%%%%%%%%%%%%%%%%%%%%%%
\subsection{Worked Example: Degree on the Circle}
\label{subsec:worked-degree-circle}
%%%%%%%%%%%%%%%%%%%%%%%%%%%%%%%%%%%%%%%%%%%%%%%%%%%%%%%%%%%%

\begin{workedexamplebox}{Degree of a Power Map}

Consider

\[
f:S^1\to S^1
\]

defined using complex notation by

\[
f(z)=z^n.
\]

As \(z\) travels once around \(S^1\), the image travels around the target
circle \(n\) signed times.

Therefore

\begin{answerbox}
\[
\boxed{
\deg(f)=n.
}
\]
\end{answerbox}

This agrees with the winding-number interpretation from Algebraic Topology.
\end{workedexamplebox}

%%%%%%%%%%%%%%%%%%%%%%%%%%%%%%%%%%%%%%%%%%%%%%%%%%%%%%%%%%%%
\subsection{Common Errors}
\label{subsec:differential-topology-common-errors}
%%%%%%%%%%%%%%%%%%%%%%%%%%%%%%%%%%%%%%%%%%%%%%%%%%%%%%%%%%%%

\subsubsection{Confusing Differential Topology with Differential Geometry}

Differential topology studies smooth structure and smooth maps.

Differential geometry additionally introduces geometric structures such as
metrics, connections, and curvature.

%%%%%%%%%%%%%%%%%%%%%%%%%%%%%%%%%%%%%%%%%%%%%%%%%%%%%%%%%%%%
\subsubsection{Confusing Homeomorphism and Diffeomorphism}

A diffeomorphism is automatically a homeomorphism.

A homeomorphism need not preserve smooth structure.

%%%%%%%%%%%%%%%%%%%%%%%%%%%%%%%%%%%%%%%%%%%%%%%%%%%%%%%%%%%%
\subsubsection{Confusing Immersion and Embedding}

An immersion may have self-intersections.

An embedding cannot identify distinct points.

%%%%%%%%%%%%%%%%%%%%%%%%%%%%%%%%%%%%%%%%%%%%%%%%%%%%%%%%%%%%
\subsubsection{Assuming an Injective Immersion Is Automatically an Embedding}

Additional topological conditions are required.

An embedding must be a homeomorphism onto its image.

%%%%%%%%%%%%%%%%%%%%%%%%%%%%%%%%%%%%%%%%%%%%%%%%%%%%%%%%%%%%
\subsubsection{Confusing Critical Points and Critical Values}

A critical point lies in the domain.

A critical value lies in the codomain.

%%%%%%%%%%%%%%%%%%%%%%%%%%%%%%%%%%%%%%%%%%%%%%%%%%%%%%%%%%%%
\subsubsection{Forgetting the Empty Preimage Case}

If

\[
f^{-1}(y)=\varnothing,
\]

then \(y\) is a regular value by definition.

%%%%%%%%%%%%%%%%%%%%%%%%%%%%%%%%%%%%%%%%%%%%%%%%%%%%%%%%%%%%
\subsubsection{Applying the Regular Value Theorem Without Checking Rank}

One must verify that every point in

\[
f^{-1}(y)
\]

is regular.

%%%%%%%%%%%%%%%%%%%%%%%%%%%%%%%%%%%%%%%%%%%%%%%%%%%%%%%%%%%%
\subsubsection{Assuming Every Level Set Is a Manifold}

A level set at a critical value may contain singularities.

Regular values guarantee smooth level sets.

%%%%%%%%%%%%%%%%%%%%%%%%%%%%%%%%%%%%%%%%%%%%%%%%%%%%%%%%%%%%
\subsubsection{Confusing Transverse with Perpendicular}

Transversality does not mean that tangent spaces meet at right angles.

It means that together they span the ambient tangent space:

\[
T_pA+T_pB=T_pM.
\]

%%%%%%%%%%%%%%%%%%%%%%%%%%%%%%%%%%%%%%%%%%%%%%%%%%%%%%%%%%%%
\subsubsection{Forgetting Dimension Conditions}

For transverse submanifolds,

\[
\dim(A\cap B)
=
\dim A+\dim B-\dim M.
\]

If the expected dimension is negative, generic transverse intersections are
empty.

%%%%%%%%%%%%%%%%%%%%%%%%%%%%%%%%%%%%%%%%%%%%%%%%%%%%%%%%%%%%
\subsubsection{Assuming Sard's Theorem Says There Are No Critical Values}

Critical values may exist.

Sard's Theorem says that their set has measure zero.

%%%%%%%%%%%%%%%%%%%%%%%%%%%%%%%%%%%%%%%%%%%%%%%%%%%%%%%%%%%%
\subsubsection{Confusing Homotopy and Isotopy}

A homotopy may pass through nonembeddings.

An isotopy must remain an embedding throughout the deformation.

%%%%%%%%%%%%%%%%%%%%%%%%%%%%%%%%%%%%%%%%%%%%%%%%%%%%%%%%%%%%
\subsubsection{Assuming Every Smooth Function Is Morse}

A smooth function may have degenerate critical points.

A Morse function requires every critical point to be nondegenerate.

%%%%%%%%%%%%%%%%%%%%%%%%%%%%%%%%%%%%%%%%%%%%%%%%%%%%%%%%%%%%
\subsubsection{Confusing Morse Index with Vector-Field Index}

The Morse index counts negative Hessian directions.

The index of an isolated vector-field zero is a degree-type invariant.

They are related in important settings but are different definitions.

%%%%%%%%%%%%%%%%%%%%%%%%%%%%%%%%%%%%%%%%%%%%%%%%%%%%%%%%%%%%
\subsubsection{Assuming Morse Index Counts Critical Points}

The Morse index classifies an individual nondegenerate critical point.

The number of critical points of index \(k\) is usually denoted \(m_k\).

%%%%%%%%%%%%%%%%%%%%%%%%%%%%%%%%%%%%%%%%%%%%%%%%%%%%%%%%%%%%
\subsubsection{Ignoring Compactness Hypotheses}

Several global results in differential topology require compactness or
properness assumptions.

Local theorems should not automatically be interpreted as global ones.

%%%%%%%%%%%%%%%%%%%%%%%%%%%%%%%%%%%%%%%%%%%%%%%%%%%%%%%%%%%%
\subsubsection{Assuming Topological Manifolds Have Unique Smooth Structures}

Smooth structure can contain information invisible to ordinary topology.

Exotic spheres and exotic \(\R^4\)'s demonstrate this phenomenon.

%%%%%%%%%%%%%%%%%%%%%%%%%%%%%%%%%%%%%%%%%%%%%%%%%%%%%%%%%%%%
\subsection{Applications and Connections}
\label{subsec:differential-topology-applications}
%%%%%%%%%%%%%%%%%%%%%%%%%%%%%%%%%%%%%%%%%%%%%%%%%%%%%%%%%%%%

\begin{applicationbox}
\textbf{Algebraic Topology.}
Degree, Euler characteristic, homology, and intersection numbers connect
smooth behavior with algebraic invariants.

\medskip

\textbf{Differential Geometry.}
Differential topology provides the underlying smooth-manifold language used
for tangent spaces, differential forms, metrics, curvature, and geodesics.

\medskip

\textbf{Topology of Manifolds.}
Regular values, transversality, Morse theory, and cobordism provide tools for
classifying and decomposing manifolds.

\medskip

\textbf{Knot Theory.}
Knots are embeddings of circles, and knot equivalence is expressed through
ambient isotopy.

\medskip

\textbf{Dynamical Systems.}
Vector fields, equilibria, indices, and global topological obstructions link
smooth dynamics with topology.

\medskip

\textbf{Optimization.}
Morse theory connects critical points and Hessians with the global topology
of optimization domains.

\medskip

\textbf{Physics.}
Smooth manifolds and vector bundles provide the mathematical setting for
many classical and modern physical theories.

\medskip

\textbf{Data Analysis.}
Manifold learning assumes high-dimensional observations may concentrate near
lower-dimensional geometric structures, motivating connections between
smooth manifolds, geometry, topology, and computation.
\end{applicationbox}

%%%%%%%%%%%%%%%%%%%%%%%%%%%%%%%%%%%%%%%%%%%%%%%%%%%%%%%%%%%%
\subsection{Exercises}
\label{subsec:differential-topology-exercises}
%%%%%%%%%%%%%%%%%%%%%%%%%%%%%%%%%%%%%%%%%%%%%%%%%%%%%%%%%%%%

\begin{exercise}[1]
What is the primary difference between differential topology and
differential geometry?
\end{exercise}

\begin{exercise}[1]
Define a smooth manifold.
\end{exercise}

\begin{exercise}[1]
Define a diffeomorphism.
\end{exercise}

\begin{exercise}[1]
Define the differential

\[
df_p.
\]
\end{exercise}

\begin{exercise}[1]
Define an immersion.
\end{exercise}

\begin{exercise}[1]
Define a submersion.
\end{exercise}

\begin{exercise}[1]
Define a smooth embedding.
\end{exercise}

\begin{exercise}[1]
What is the difference between an immersion and an embedding?
\end{exercise}

\begin{exercise}[1]
Define a critical point.
\end{exercise}

\begin{exercise}[1]
Define a regular value.
\end{exercise}

\begin{exercise}[1]
State the Regular Value Theorem.
\end{exercise}

\begin{exercise}[1]
Define codimension.
\end{exercise}

\begin{exercise}[1]
State Sard's Theorem.
\end{exercise}

\begin{exercise}[1]
Define transverse intersection.
\end{exercise}

\begin{exercise}[1]
What does

\[
A\pitchfork B
\]

mean?
\end{exercise}

\begin{exercise}[1]
Define an orientable manifold.
\end{exercise}

\begin{exercise}[1]
Define a manifold with boundary.
\end{exercise}

\begin{exercise}[1]
Define an isotopy.
\end{exercise}

\begin{exercise}[1]
Define a Morse function.
\end{exercise}

\begin{exercise}[1]
Define a nondegenerate critical point.
\end{exercise}

\begin{exercise}[1]
Define the Morse index.
\end{exercise}

\begin{exercise}[1]
State the Morse Lemma.
\end{exercise}

\begin{exercise}[1]
What is an \(n\)-dimensional handle of index \(\lambda\)?
\end{exercise}

\begin{exercise}[1]
Define cobordism.
\end{exercise}

\begin{exercise}[2]
Show that

\[
S^1
=
\{(x,y)\in\R^2:x^2+y^2=1\}
\]

is a smooth submanifold of \(\R^2\).
\end{exercise}

\begin{exercise}[2]
Show that

\[
S^2
=
\{(x,y,z)\in\R^3:x^2+y^2+z^2=1\}
\]

is a smooth submanifold of \(\R^3\).
\end{exercise}

\begin{exercise}[2]
Determine whether \(0\) is a regular value of

\[
f(x,y)=x^2+y^2.
\]
\end{exercise}

\begin{exercise}[2]
Determine whether \(1\) is a regular value of

\[
f(x,y)=x^2+y^2.
\]
\end{exercise}

\begin{exercise}[2]
Find the critical points of

\[
f(x,y)=x^2+y^2.
\]
\end{exercise}

\begin{exercise}[2]
Find the critical points of

\[
f(x,y)=x^2-y^2.
\]
\end{exercise}

\begin{exercise}[2]
Find the Morse index of the critical point of

\[
f(x,y)=x^2-y^2.
\]
\end{exercise}

\begin{exercise}[2]
Classify the critical point of

\[
f(x,y)=x^2+y^2
\]

using the Hessian.
\end{exercise}

\begin{exercise}[2]
Classify the critical point of

\[
f(x,y)=-x^2-y^2.
\]
\end{exercise}

\begin{exercise}[2]
Explain why a transverse intersection of two surfaces in a
three-dimensional manifold is generically one-dimensional.
\end{exercise}

\begin{exercise}[2]
If

\[
A^3,B^4\subseteq M^6
\]

intersect transversely, determine

\[
\dim(A\cap B).
\]
\end{exercise}

\begin{exercise}[2]
Explain why perpendicular submanifolds are transverse but transverse
submanifolds need not be perpendicular.
\end{exercise}

\begin{exercise}[2]
Find the degree of

\[
f:S^1\to S^1,
\qquad
f(z)=z^5.
\]
\end{exercise}

\begin{exercise}[2]
Find the degree of

\[
f:S^1\to S^1,
\qquad
f(z)=z^{-3}.
\]
\end{exercise}

\begin{exercise}[2]
Explain why a map of nonzero degree between compact connected oriented
manifolds of equal dimension must be surjective.
\end{exercise}

\begin{exercise}[2]
Explain why the Hairy Ball Theorem applies to \(S^2\).
\end{exercise}

\begin{exercise}[2]
Construct a nowhere-zero tangent vector field on \(S^1\).
\end{exercise}

\begin{exercise}[2]
Verify

\[
\chi(S^2)=2
\]

and explain its relationship with the Poincar\'e--Hopf Theorem.
\end{exercise}

\begin{exercise}[2]
For a standard height function on \(S^2\), list the critical points and
their Morse indices.
\end{exercise}

\begin{exercise}[2]
For a standard height function on a torus, list the four critical-point
indices.
\end{exercise}

\begin{exercise}[2]
Use the Morse critical points of the torus to calculate its Euler
characteristic.
\end{exercise}

\begin{exercise}[2]
Explain why isotopy is stronger than homotopy.
\end{exercise}

\begin{exercise}[2]
Give an example of two maps that may be homotopic even though a proposed
deformation through embeddings would be impossible.
\end{exercise}

\begin{exercise}[3]
Prove that a diffeomorphism is a homeomorphism.
\end{exercise}

\begin{exercise}[3]
Show that the inclusion

\[
S^n\hookrightarrow\R^{n+1}
\]

is an immersion.
\end{exercise}

\begin{exercise}[3]
Show that the projection

\[
\pi:\R^{m+n}\to\R^n
\]

onto the final \(n\) coordinates is a submersion.
\end{exercise}

\begin{exercise}[3]
Use the Regular Value Theorem to show that

\[
\{(x,y,z)\in\R^3:x^2+y^2-z^2=1\}
\]

is a smooth surface.
\end{exercise}

\begin{exercise}[3]
Find all critical values of

\[
f(x,y)=x^2+y^2.
\]
\end{exercise}

\begin{exercise}[3]
For

\[
f(x,y)=x^3-3x+y^2,
\]

find the critical points and determine whether each is nondegenerate.
\end{exercise}

\begin{exercise}[3]
Determine the Morse index of each nondegenerate critical point in the
previous exercise.
\end{exercise}

\begin{exercise}[3]
Explain how the Constant Rank Theorem contains the local behavior of
immersions and submersions as special cases.
\end{exercise}

\begin{exercise}[3]
Prove that if \(f:M\to N\) is transverse to a point \(y\in N\), then \(y\)
is a regular value.
\end{exercise}

\begin{exercise}[3]
Suppose

\[
A^a,B^b\subseteq M^n
\]

intersect transversely. Derive the dimension formula

\[
\dim(A\cap B)=a+b-n.
\]
\end{exercise}

\begin{exercise}[3]
Explain why generic perturbations are useful when defining intersection
numbers.
\end{exercise}

\begin{exercise}[3]
Show that

\[
\partial(\partial M)=\varnothing
\]

for familiar examples of manifolds with boundary.
\end{exercise}

\begin{exercise}[3]
Compute the degree of the antipodal map

\[
A:S^n\to S^n,
\qquad
A(x)=-x.
\]
\end{exercise}

\begin{exercise}[3]
Show that the degree of the antipodal map is

\[
(-1)^{n+1}.
\]
\end{exercise}

\begin{exercise}[3]
Explain why degree is invariant under homotopy.
\end{exercise}

\begin{exercise}[3]
Use the Poincar\'e--Hopf theorem to explain why a compact manifold with
nonzero Euler characteristic cannot admit a nowhere-zero vector field.
\end{exercise}

\begin{exercise}[3]
For a Morse function on a compact surface, explain how index-\(0\),
index-\(1\), and index-\(2\) critical points alter the topology of sublevel
sets.
\end{exercise}

\begin{exercise}[3]
Verify the weak Morse inequalities for \(S^2\) using its standard height
function.
\end{exercise}

\begin{exercise}[3]
Verify the weak Morse inequalities for \(T^2\) using a generic height
function.
\end{exercise}

\begin{exercise}[3]
Explain how Morse theory gives a handle decomposition of a compact smooth
manifold.
\end{exercise}

\begin{exercise}[3]
Show that

\[
S^1\times[0,1]
\]

is a cobordism between two circles.
\end{exercise}

\begin{exercise}[4]
Study the proof of the Whitney Embedding Theorem and explain why the
dimension \(2n\) appears.
\end{exercise}

\begin{exercise}[4]
Study Sard's Theorem and explain why it is essential in degree theory.
\end{exercise}

\begin{exercise}[4]
Research the Parametric Transversality Theorem.
\end{exercise}

\begin{exercise}[4]
Use transversality to explain why two generic curves in a surface intersect
in isolated points.
\end{exercise}

\begin{exercise}[4]
Use transversality to explain why two generic surfaces in a
three-dimensional manifold intersect along curves.
\end{exercise}

\begin{exercise}[4]
Research the relationship between intersection number and homology.
\end{exercise}

\begin{exercise}[4]
Research one differential-topological proof of the Brouwer Fixed-Point
Theorem.
\end{exercise}

\begin{exercise}[4]
Research the relationship between degree theory and the Fundamental Theorem
of Algebra.
\end{exercise}

\begin{exercise}[4]
Study the proof of the Hairy Ball Theorem using degree or homology.
\end{exercise}

\begin{exercise}[4]
Research the proof of the Poincar\'e--Hopf Theorem.
\end{exercise}

\begin{exercise}[4]
Construct a Morse function on the torus and analyze its critical points.
\end{exercise}

\begin{exercise}[4]
Research the strong Morse inequalities.
\end{exercise}

\begin{exercise}[4]
Explain how a Morse function gives a CW decomposition of a manifold.
\end{exercise}

\begin{exercise}[4]
Research Morse homology and describe how its boundary operator is defined.
\end{exercise}

\begin{exercise}[4]
Research exotic spheres and explain the distinction between homeomorphism
and diffeomorphism that they demonstrate.
\end{exercise}

\begin{exercise}[4]
Research exotic \(\R^4\) and explain why dimension four is exceptional in
smooth topology.
\end{exercise}

\begin{exercise}[4]
Investigate cobordism groups and explain how disjoint union defines their
addition.
\end{exercise}

%%%%%%%%%%%%%%%%%%%%%%%%%%%%%%%%%%%%%%%%%%%%%%%%%%%%%%%%%%%%
\subsection{Conceptual Summary}
\label{subsec:differential-topology-summary}
%%%%%%%%%%%%%%%%%%%%%%%%%%%%%%%%%%%%%%%%%%%%%%%%%%%%%%%%%%%%

Differential topology studies smooth manifolds and smooth maps while
ignoring metric quantities such as length and curvature.

The differential

\[
\boxed{
df_p:T_pM\to T_{f(p)}N
}
\]

encodes the local linear behavior of a smooth map.

An immersion has injective differential.

A submersion has surjective differential.

An embedding is an immersion that is also a homeomorphism onto its image.

Regular values generate smooth submanifolds:

\[
\boxed{
\dim f^{-1}(y)
=
\dim M-\dim N.
}
\]

Sard's Theorem shows that critical values form a measure-zero set.

Transversality generalizes regularity and makes intersections behave with
the expected dimension:

\[
\boxed{
\dim(A\cap B)
=
\dim A+\dim B-\dim M.
}
\]

Orientation allows intersections and preimages to be counted with signs.

For maps between compact connected oriented manifolds of equal dimension,
these signs produce the degree:

\[
\boxed{
\deg(f)
=
\sum_{x\in f^{-1}(y)}
\operatorname{sgn}(df_x).
}
\]

Vector fields connect local differential behavior with global topology
through

\[
\boxed{
\sum_p\operatorname{ind}_p(X)
=
\chi(M).
}
\]

Morse theory studies manifolds using critical points of functions

\[
f:M\to\R.
\]

At a nondegenerate critical point,

\[
\boxed{
f
=
f(p)
-
x_1^2-\cdots-x_\lambda^2
+
x_{\lambda+1}^2+\cdots+x_n^2
}
\]

in suitable local coordinates.

Passing a critical value corresponds to attaching a handle

\[
\boxed{
D^\lambda\times D^{n-\lambda}.
}
\]

The numbers of Morse critical points are constrained by homology through the
Morse inequalities.

Cobordism studies manifolds that arise together as boundaries of
higher-dimensional manifolds.

Finally, exotic smooth structures demonstrate that smooth topology can
contain information invisible to ordinary topology.

Thus differential topology creates the bridge

\[
\boxed{
\text{calculus}
\longleftrightarrow
\text{smooth manifolds}
\longleftrightarrow
\text{global topology}.
}
\]

%%%%%%%%%%%%%%%%%%%%%%%%%%%%%%%%%%%%%%%%%%%%%%%%%%%%%%%%%%%%
\subsection{Section Connections}
\label{subsec:differential-topology-connections}
%%%%%%%%%%%%%%%%%%%%%%%%%%%%%%%%%%%%%%%%%%%%%%%%%%%%%%%%%%%%

\chapterconnection{
Differential Topology builds directly on Point-Set Topology and Algebraic
Topology while using the smooth-manifold language introduced in Differential
Geometry. The Regular Value Theorem and Transversality Theorem explain how
smooth submanifolds and intersections arise. Degree theory connects smooth
maps with homology and orientation. The Poincar\'e--Hopf Theorem connects
vector fields with Euler characteristic. Morse Theory converts critical-point
data into global topological information and provides handle decompositions
of manifolds. Isotopy provides the natural language for embedded objects and
leads directly into Knot Theory. Cobordism and smooth structures become
central in higher-dimensional manifold topology. These ideas prepare the
reader for Geometric Topology and the dedicated study of Manifolds that
follows.
}

%%%%%%%%%%%%%%%%%%%%%%%%%%%%%%%%%%%%%%%%%%%%%%%%%%%%%%%%%%%%
% SECTION SOURCE TRACKING
%%%%%%%%%%%%%%%%%%%%%%%%%%%%%%%%%%%%%%%%%%%%%%%%%%%%%%%%%%%%

%------------------------------------------------------------
% 4.3 Differential Topology
%------------------------------------------------------------
%
% Source basis:
%   - Standard differential topology and smooth manifold
%     theory.
%
% Used for:
%   - smooth manifolds;
%   - smooth maps and diffeomorphisms;
%   - tangent spaces and differentials;
%   - rank;
%   - immersions;
%   - submersions;
%   - embeddings;
%   - Whitney embedding and immersion theorems;
%   - critical and regular points;
%   - critical and regular values;
%   - Regular Value Theorem;
%   - level sets and codimension;
%   - Constant Rank Theorem;
%   - inverse and implicit function principles;
%   - Sard's Theorem;
%   - transversality;
%   - Preimage Theorem;
%   - genericity;
%   - intersection numbers;
%   - orientations;
%   - manifolds with boundary;
%   - degree and mod-2 degree;
%   - vector fields and zeros;
%   - Hairy Ball Theorem;
%   - vector-field index;
%   - Poincare--Hopf Theorem;
%   - smooth homotopy;
%   - isotopy and ambient isotopy;
%   - Morse theory;
%   - Hessians and nondegenerate critical points;
%   - Morse index;
%   - Morse Lemma;
%   - sublevel sets;
%   - handles;
%   - Morse inequalities;
%   - Morse homology preview;
%   - cobordism;
%   - smooth structures;
%   - exotic spheres;
%   - exotic R^4;
%   - applications to knots, dynamics,
%     optimization, and physics.
%
% Notes:
%   - General manifold definitions are developed further in
%     the dedicated Manifold Theory section.
%   - Tangent spaces, differential forms, metrics, curvature,
%     and connections are developed canonically in
%     Differential Geometry.
%   - Homology, cohomology, degree, and Euler characteristic
%     are introduced in Algebraic Topology.
%   - Knot embeddings and ambient isotopy are developed in
%     the Knot Theory section.
%   - Morse homology is introduced only conceptually here.
%   - Cobordism theory is introduced at structural level;
%     characteristic classes and advanced classification
%     belong to later manifold/topology material.
%
%%%%%%%%%%%%%%%%%%%%%%%%%%%%%%%%%%%%%%%%%%%%%%%%%%%%%%%%%%%%